\begin{document}
\maketitle

\begin{abstract}
We isolate the uniformly finite-dimensional, finite-range Hamiltonian class
actually assembled by the TFPT \(2+1\)- and \(3+1\)-dimensional probes:
electric and plaquette gauge terms, gauge-dressed fermion hopping, the
DET16 mirror projector, and the seam clock and charge terms.  At the frozen
\(2+1\)-dimensional coefficients the interaction has range \(R=2\), largest
local term \(J=12/5\), and interaction moment
\(\kappa=16043/450\).  A complete finite-range Lieb--Robinson argument gives
\[
 v_{\rm LR}=2e\kappa R=\frac{32086e}{225}
 =387.639069990831.
\]
We then prove in norm, rather than assume, that the finite-volume
Heisenberg dynamics form a Cauchy net on every local observable and hence
define a strongly continuous thermodynamic-limit automorphism group.
Gauss constraints are handled on the gauge-invariant subalgebra of the
ambient quasilocal algebra; naive embeddings of projected physical Hilbert
spaces are neither needed nor asserted.  Weak-* compactness gives at least
one ground state and, for each positive inverse temperature, at least one
KMS state.  The numeric twin passes all eight checks and measures velocity
\(2.0\), well below the proved class bound.
\end{abstract}

\tableofcontents

\section{The concrete uniformly local class}

\subsection{Ambient tensor-product algebra}

Let \(\Gamma=\Z^2\) with graph distance \(d\) and coordination \(z=4\).
For a finite \(\Lambda\Subset\Gamma\), put
\[
 \Hh_\Lambda=\bigotimes_{x\in\Lambda}\Hh_x,\qquad
 \A_\Lambda=\mathcal B(\Hh_\Lambda),
 \qquad
 \iota_{\Lambda,\Lambda'}(A)=A\otimes\id_{\Lambda'\setminus\Lambda}.
\tag{1.1}\label{eq:ambient}
\]
One cell contains the four representative physical edge modes, a mirror
Fock factor, the two positively oriented outgoing finite-group gauge links,
and a four-state seam clock.  Thus
\[
 \dim\Hh_x=16\cdot 2^{n_{\rm mir}}\cdot 2^2\cdot4,
 \qquad n_{\rm mir}\in\{4,16\}.
\tag{1.2}\label{eq:localdim}
\]
The DET4 numerical analogue has dimension \(4096\); the full DET16 cell has
dimension \(16\,777\,216\).  Only finiteness and a volume-independent upper
bound are used below.  The quasilocal algebra is
\[
 \A=\overline{\bigcup_{\Lambda\Subset\Gamma}\A_\Lambda}^{\norm{\cdot}}.
\tag{1.3}\label{eq:quasilocal}
\]

\begin{definition}[TFPT assembled interaction]
For finite \(X\Subset\Gamma\), let \(\Phi(X)=\Phi(X)^*\), with
\(\Phi(X)=0\) unless \(X\) is one of the supports in
Table~\ref{tab:terms}.  The free-boundary finite-volume Hamiltonian is
\[
 H_\Lambda=\sum_{X\subseteq\Lambda}\Phi(X).
\tag{1.4}\label{eq:H}
\]
Scalar multiples of the identity may be removed termwise, since they do
not change \(e^{itH_\Lambda}Ae^{-itH_\Lambda}\).
\end{definition}

\begin{table}[ht]
\centering
\small
\begin{tabularx}{\textwidth}{@{}lXccc@{}}
\toprule
term & frozen operator and coefficient & \(\norm{\Phi(X)}\) & \(|X|\) &
\(\diam X\)\\
\midrule
electric &
\(\frac{g^2}{2}E_\ell^2,\ g=1.2\), with each positive link assigned to
its tail cell & \(18/25=0.72\) & 1 & 0\\
plaquette &
\(\frac1{2g^2}(2-U_p-U_p^*)\); its dynamical representative is
\(-\frac1{2g^2}(U_p+U_p^*)\) & \(25/36\) & 4 & 2\\
hopping &
\(-t\sum_{a=1}^4(c_{y,a}^*U_{xy}c_{x,a}+\mathrm{h.c.}),\ t=0.6\)
& \(\leq4t=12/5\) & 2 & 1\\
DET16 &
\(2(\id-P_{\phi,x})\), where \(P_{\phi,x}\) is the rank-one
Spin(10)-invariant determinant projector & \(2\) & 1 & 0\\
seam clock &
\(-0.5(\Psi_x+\Psi_x^*)\), \(\Psi_x^4=\id\) & \(1\) & 1 & 0\\
seam charge &
\(0.2(\Psi_x+\Psi_x^*) (n_{Q,x}-1/2)/2\) & \(1/10\) & 1 & 0\\
staggered mass &
\(0.2\sum_{a=1}^4(-1)^x(n_{x,a}-n^{\rm bg}_{x,a})\) & \(\leq4/5\) & 1 & 0\\
\bottomrule
\end{tabularx}
\caption{Supports and norm bounds extracted from the frozen assembled
models.  The positive plaquette term itself has norm at most
\(2/g^2=25/18\); after deleting its scalar part, the norm relevant to
dynamics and commutators is \(1/g^2=25/36\).}
\label{tab:terms}
\end{table}

The determinant identity checked in
\texttt{verification/v1002\_spin10\_mirror\_projector.py} gives
\(P_{\phi,x}=P_{\phi,x}^*=P_{\phi,x}^2\), so its row has norm exactly
two.  The clock and charge bounds follow from
\(\norm{\Psi}=1\), \(\norm{(\Psi+\Psi^*)/2}=1\), and
\(\norm{n_Q-1/2}=1/2\).

\subsection{Uniform constants}

Define the finite-range interaction moment
\[
 \kappa:=\sup_{x\in\Gamma}
 \sum_{X\ni x}|X|\,\norm{\Phi(X)}.
\tag{1.5}\label{eq:kappa}
\]
There are two assigned electric links, four plaquettes, and four bonds
incident on a square-lattice cell.  Consequently
\[
\begin{aligned}
 J&:=\sup_X\norm{\Phi(X)}=\frac{12}{5},\qquad R=2,\\
 \kappa
 &\leq
 2\frac{18}{25}
 +4\cdot4\frac{25}{36}
 +4\cdot2\frac{12}{5}
 +2+1+\frac1{10}+\frac45\\
 &=\frac{16043}{450}=35.651111111111\ldots .
\end{aligned}
\tag{1.6}\label{eq:constants}
\]
These bounds are independent of \(\Lambda\).  More generally, bounded
degree, finite \(R\), finite local dimension and finite \(\kappa\) are the
class hypotheses; the displayed numbers are the frozen assembled
\(2+1\)-dimensional specialization.

\section{The Lieb--Robinson theorem and explicit TFPT velocity}

\begin{theorem}[Finite-range Lieb--Robinson theorem, full form used here]
\label{thm:LRclassical}
Let \((\Gamma,d)\) be a graph, let every local Hilbert space be finite
dimensional with a uniform dimension bound, and let a self-adjoint
interaction satisfy
\[
 \Phi(X)=0\quad\text{if }\diam X>R,\qquad
 \sup_x\sum_{X\ni x}|X|\norm{\Phi(X)}=\kappa<\infty.
\tag{2.1}\label{eq:LRhyp}
\]
For every finite \(\Lambda\), define
\(\tau_t^\Lambda(A)=e^{itH_\Lambda}Ae^{-itH_\Lambda}\).
If \(A\in\A_X\), \(B\in\A_Y\), and \(d(X,Y)>0\), then, uniformly in
\(\Lambda\supset X\cup Y\),
\[
 \norm{[\tau_t^\Lambda(A),B]}
 \leq
 2\norm A\norm B\,|X|\,
 \exp\!\left[-\frac{d(X,Y)}R+2e\kappa|t|\right].
\tag{2.2}\label{eq:LR}
\]
Thus an admissible Lieb--Robinson velocity is
\[
 v_{\rm LR}=2e\kappa R.
\tag{2.3}\label{eq:vLR}
\]
The same assertion holds for even observables in a CAR algebra, because
disjoint even local algebras commute.
\end{theorem}

\begin{proof}
This is the original Lieb--Robinson iteration in the finite-range
normalization \eqref{eq:LRhyp}; it is equivalent to the exponential
Nachtergaele--Sims form.  We include the complete estimate.
Differentiate the commutator after removing all Hamiltonian terms supported
entirely away from its evolving support.  Duhamel iteration gives
\[
 \norm{[\tau_t^\Lambda(A),B]}
 \leq2\norm A\norm B
 \sum_{n\geq1}\frac{(2|t|)^n}{n!}\,a_n(X,Y),
\tag{2.4}\label{eq:iteration}
\]
where \(a_n(X,Y)\) is the sum of
\(\prod_{j=1}^n\norm{\Phi(Z_j)}\) over chains with
\(Z_1\cap X\ne\varnothing\), \(Z_j\cap Z_{j-1}\ne\varnothing\), and
\(Z_n\cap Y\ne\varnothing\).  Choosing one intersection site at each
step and using \eqref{eq:LRhyp} yields
\[
 a_n(X,Y)\leq |X|\kappa^n.
\tag{2.5}\label{eq:pathweight}
\]
Every member of a chain has diameter at most \(R\); hence no such chain can
join \(X\) to \(Y\) unless \(nR\geq d(X,Y)\).  Put
\(m=\lceil d(X,Y)/R\rceil\).  For \(n\geq m\),
\(\mathbf1_{\{n\geq m\}}\leq e^{n-m}\), and therefore
\[
 \sum_{n\geq m}\frac{(2\kappa|t|)^n}{n!}
 \leq e^{-m}\sum_{n\geq0}\frac{(2e\kappa|t|)^n}{n!}
 \leq
 e^{-d(X,Y)/R+2e\kappa|t|}.
\tag{2.6}\label{eq:poissontail}
\]
Substitution in \eqref{eq:iteration} proves \eqref{eq:LR}.  Rewriting
the exponent as
\(-[d(X,Y)-2e\kappa R|t|]/R\) proves \eqref{eq:vLR}.  In the fermionic
case all Hamiltonian terms and physical observables under consideration
are even, so the initial disjoint-support commutator and every step of the
iteration are ordinary commutators.
\end{proof}

The theorem is the finite-range specialization of Lieb and
Robinson~\cite{LR1972} and of the exponential bounds of
Nachtergaele and Sims~\cite{NS2006,NS2010}.  For the TFPT constants,
\[
 \boxed{
 v_{\rm LR}=2e\frac{16043}{450}\,2
 =\frac{32086e}{225}
 =387.639069990831.}
\tag{2.7}\label{eq:TFPTvelocity}
\]
This deliberately conservative class bound includes all four flavors,
plaquettes, and onsite terms.  It is not a fitted physical signal speed.

\section{Norm thermodynamic limit}

\begin{theorem}[Thermodynamic-limit dynamics]
\label{thm:dynamics}
For every \(A\in\A_{\rm loc}:=\bigcup_{\Lambda}\A_\Lambda\), the net
\[
 \tau_t(A):=\lim_{\Lambda\uparrow\Gamma}\tau_t^\Lambda(A)
\tag{3.1}\label{eq:limit}
\]
exists in norm, uniformly for \(t\) in compact real intervals and
independently of the exhausting sequence.  It extends uniquely to a
strongly continuous one-parameter group of \(C^*\)-automorphisms of
\(\A\).
\end{theorem}

\begin{proof}
We give the standard Bratteli--Robinson/Nachtergaele--Sims proof with an
explicit Cauchy bound.  Let \(A\in\A_X\), and let
\(\Lambda\subset\Lambda'\) contain the closed \(b\)-neighbourhood of \(X\).
Write \(V=H_{\Lambda'}-H_\Lambda\), with both Hamiltonians embedded in
\(\A_{\Lambda'}\).  Differentiating
\(\tau_{t-s}^{\Lambda'}(\tau_s^\Lambda(A))\) gives
\[
 \tau_t^{\Lambda'}(A)-\tau_t^\Lambda(A)
 =i\int_0^t
 \tau_{t-s}^{\Lambda'}
 \bigl([V,\tau_s^\Lambda(A)]\bigr)\,ds.
\tag{3.2}\label{eq:duhamel}
\]
Every term \(Z\) in \(V\) reaches outside \(\Lambda\).  Since
\(\diam Z\leq R\), \(d(X,Z)\geq\rho:=\max(0,b-R)\).
Apply \eqref{eq:LR} termwise.  If \(q=e^{-1/R}\), then the square-lattice
shell bound
\[
 \sum_{y:d(X,y)\geq\rho}e^{-d(X,y)/R}
 \leq
 4|X|q^\rho
 \left(\frac{\rho+1}{1-q}+\frac{q}{(1-q)^2}\right)
 =:|X|S_R(\rho)
\tag{3.3}\label{eq:shell}
\]
follows by assigning each \(y\) to a nearest point of \(X\) and using at
most \(4(n+1)\) sites in shell \(n\).  Assign each interaction \(Z\) to a
site \(y\in Z\) nearest \(X\).  Because \(\diam Z\leq R\),
\(e^{-d(X,Z)/R}\leq e\,e^{-d(X,y)/R}\), while
\(\sum_{Z\ni y}\norm{\Phi(Z)}\leq\kappa\).  Equations
\eqref{eq:duhamel}--\eqref{eq:shell} therefore give, for \(t\geq0\),
\[
\begin{aligned}
 \norm{\tau_t^{\Lambda'}(A)-\tau_t^\Lambda(A)}
 &\leq
 2\norm A\,|X|\,e\kappa |X|S_R(\rho)
 \int_0^t e^{2e\kappa s}\,ds\\
 &=\norm A\,|X|^2S_R(\rho)
 \bigl(e^{2e\kappa t}-1\bigr).
\end{aligned}
\tag{3.4}\label{eq:convergence}
\]
The same bound holds with \(|t|\).  For fixed \(T\), its right-hand side
tends to zero as \(b\to\infty\), uniformly for \(|t|\leq T\).  Hence the
finite-volume dynamics are a norm-Cauchy net on local observables; this
also proves independence of the exhaustion.

Each \(\tau_t^\Lambda\) is isometric, linear, multiplicative and
star-preserving.  Norm limits preserve these properties.
Applying the construction to \(-t\) shows that \(\tau_{-t}\) is the
inverse, and taking limits in
\(\tau_{s+t}^\Lambda=\tau_s^\Lambda\tau_t^\Lambda\) gives the group law.
For local \(A\), continuity at \(t=0\) follows first in a fixed sufficiently
large volume and then from the compact-time uniform limit.  Density of
\(\A_{\rm loc}\) and isometry extend both \(\tau_t\) and strong continuity
to all of \(\A\).  This is precisely the conclusion of the standard
thermodynamic-dynamics theorem, commonly cited as
Bratteli--Robinson, Theorem~6.2.11~\cite{BR}.
\end{proof}

\section{Gauss law and the correct consistent family}

Let \(\mathcal G_{\rm loc}\) be the group of finitely supported local gauge
transformations on the ambient field algebra, and let \(\alpha_g\) denote
its action.  Define
\[
 \A^G=\{A\in\A:\alpha_g(A)=A\text{ for every }g\in\mathcal G_{\rm loc}\}.
\tag{4.1}\label{eq:fixed}
\]
The electric and plaquette terms are gauge invariant; the hopping contains
its parallel transporter; the determinant projector is a Spin(10) singlet;
and the seam charge is neutral as assembled.  Thus
\[
 \alpha_g(\Phi(X))=\Phi(X),\qquad
 \alpha_g\circ\tau_t^\Lambda=\tau_t^\Lambda\circ\alpha_g.
\tag{4.2}\label{eq:gaugecov}
\]

\begin{proposition}[Compatible constrained dynamics]
\label{prop:gauss}
The ambient embeddings \eqref{eq:ambient} restrict to an inductive family
of local gauge-invariant observables, \(\tau_t(\A^G)=\A^G\), and
\(\tau_t|_{\A^G}\) is the thermodynamic dynamics of the constrained theory.
For every finite region \(F\), let \(P_F\) be the product/Haar Gauss
projection.  If \(A\) is gauge invariant and supported in the interior of
\(F\), then
\[
 [\iota(A),P_F]=0,\qquad
 P_F\iota(A)P_F=\iota(A)P_F.
\tag{4.3}\label{eq:projectorcompat}
\]
\end{proposition}

\begin{proof}
If \(A\) is fixed by every local gauge transformation, so is
\(A\otimes\id\); hence the ambient embedding preserves the fixed-point
condition.  Equation \eqref{eq:gaugecov} implies that every
\(\tau_t^\Lambda(A)\), and therefore its norm limit, is gauge invariant.
This proves invariance and consistency of the restricted automorphism
group.

For a finite gauge group,
\(P_F=\prod_{x\in F}|G|^{-1}\sum_{g\in G}U_x(g)\); for a compact group,
replace each sum by normalized Haar integration.  Gauge invariance says
\(U_x(g)\iota(A)U_x(g)^*=\iota(A)\) for every factor.  Thus \(\iota(A)\)
commutes with each averaged projector and with their product, proving
\eqref{eq:projectorcompat}.  This is the standard
gauge-invariant-subalgebra treatment of the local Gauss law
\cite{GrundlingLledo,GrundlingRudolph}.
\end{proof}

\begin{remark}[What is not compatible]
The projected spaces \(P_\Lambda\Hh_\Lambda\) do not generally embed as
\(P_\Lambda\Hh_\Lambda\otimes\Hh_{\Lambda'\setminus\Lambda}\subset
P_{\Lambda'}\Hh_{\Lambda'}\): boundary Gauss laws correlate matter with
new links and boundary flux sectors.  Therefore no naive tensor embedding
of physical Hilbert spaces is claimed.  The honest constrained
thermodynamic object is \(\A^G\), for which
\eqref{eq:projectorcompat} is exactly the required projector--restriction
commutation.  A constraint-violating operator need not commute with
\(P_F\); the numeric mutant has commutator norm one.
\end{remark}

